\begin{figure}[t]
\centering
\setlength{\fboxsep}{6pt}
\begin{minipage}[t]{0.48\textwidth}
\centering
\fbox{\begin{minipage}[t][5.2cm][t]{0.91\linewidth}
\centering
\textbf{Orbit 649 mechanism}\par\medskip
\begin{tikzpicture}[scale=.78,every node/.style={font=\small}]
\node[treev] (u1) at (-1.45,.45) {$U$};
\node[reticv] (v1) at (1.45,.45) {$V$};
\node[reticv] (x1) at (0,-1.0) {$X$};
\draw[reticedge] (u1)--(v1); \draw[reticedge] (u1)--(x1); \draw[reticedge] (v1)--(x1);
\node[port] (e1) at (0,-1.85) {$e\in\{5,6,7\}$}; \draw[ordinary] (x1)--(e1);
\end{tikzpicture}
\vspace{-2mm}
\[
\begin{aligned}
F_{12e}
={}&M_A(1-\lambda_V)(1-\lambda_X)(1-a)^2\\
&\times\bigl[\lambda_V(1-\lambda_X)b+\lambda_Xc\bigr]>0.
\end{aligned}
\]
\end{minipage}}
\end{minipage}\hfill
\begin{minipage}[t]{0.48\textwidth}
\centering
\fbox{\begin{minipage}[t][5.2cm][t]{0.91\linewidth}
\centering
\textbf{Orbit 705 mechanism}\par\medskip
\begin{tikzpicture}[scale=.78,every node/.style={font=\small}]
\node[treev] (u2) at (-1.45,.45) {$U$};
\node[reticv] (v2) at (1.45,.45) {$V$};
\node[reticv] (x2) at (0,-1.0) {$X$};
\draw[reticedge] (u2)--(v2); \draw[reticedge] (u2)--(x2); \draw[reticedge] (v2)--(x2);
\node[port] (e2) at (2.15,-.25) {$e$}; \draw[ordinary] (v2)--(e2);
\end{tikzpicture}
\vspace{2mm}
\[
F_{12e}=M_B\lambda_V(1-\lambda_V)(1-a)^2>0.
\]
\end{minipage}}
\end{minipage}
\caption{The two genuine seven-port residual mechanisms.  In both panels the competing cycle has sink label $4$, and hence $F_{12e}=0$.  The polynomial $F_{abc}$ is the established JC trinet inequality; the new step is choosing a restored support label $e$ that makes it a strict separator on the completed core-3 theta target.}
\label{fig:residual-orbits}
\end{figure}
